\section{Seguridad de la política y despliegue sistematizado}

\subsection{Risk-aware fine-tuning}

\begin{figure}[H]
\centering
\includegraphics[width=0.88\textwidth]{slide-30.jpg}
\caption{El pipeline consciente del riesgo que resume la diapositiva 30 de la clase: primero se usa el conservative critic para acotar el valor, y luego se hace un safety check antes del despliegue.}
\end{figure}

La clase menciona una capa de seguridad: antes del actor update se revisa la salida del critic, y si el valor de la política actual rebasa el pre-defined threshold, se retrocede al paso anterior. Esto es un análogo al trust region, pero está más orientado a servir de deployment guard.

\begin{warningbox}{Recomendación sobre el conservative critic}
Se guardan varios conjuntos de critic checkpoints, y al desplegar se elige el conjunto conservative. Así, aunque el newest critic overestimates, el actor tampoco será inducido a error por una señal demasiado fuerte.
\end{warningbox}

\subsection{Monitoreo en línea y rollback}

En los sistemas de producción, una práctica común es enviar los promedios de KL, critic loss y entropy al dashboard. Si cualquier métrica supera el umbral tres veces seguidas, se dispara automáticamente un rollback y se reinicia el entropy bonus.

\begin{knowledgebox}{Experiencia de los sistemas en producción}
En los sistemas de LLM o RL, la señal de early warning de policy divergence es muy importante. La clase recomienda calcular el KL con una 3-day rolling window, y si supera 2 veces la historic std, pausar la actualización del actor para que el critic catch up.
\end{knowledgebox}

\subsection{Resumen de la sección}

Un despliegue seguro exige que el sistema actor-critic forme un cortafuegos entre la actualización y la inference: conservative critic, guard rails, rollback automático y monitoreo continuo.

